\documentclass[journal]{IEEEtran}
\usepackage{cite}
\usepackage{array,booktabs,tabularx}
\usepackage{hyperref}
\hypersetup{colorlinks=true,citecolor=blue,urlcolor=blue,linkcolor=blue}
\title{Detecting Coordinated DIS Flooding and DAO Manipulation Attacks in RPL-Based IoT Networks via Adaptive Evidence Fusion}
\author{Karim Chehab\\[0.5ex]
\normalsize Department of Computer Science and Mathematics\\
\normalsize Lebanese American University, Beirut, Lebanon\\
\normalsize Email: \href{mailto:karim.chehab02@lau.edu}{karim.chehab02@lau.edu}}
\date{}
\begin{document}
\maketitle

\section{Introduction}\label{sec:introduction}
A routing defense must protect communication as well as identify malicious behavior. This requirement motivates the present study of constrained industrial Internet of Things (IIoT) networks. A compromised routing node may repeatedly solicit routing information while advertising an incorrect parent for downward routes. The solicitations add control traffic; the false advertisement can damage the route used to reach a destination. Assessing a defense against this combination therefore requires attention to both its decisions and the service available to legitimate nodes.

The IPv6 Routing Protocol for Low-Power and Lossy Networks (RPL) organizes nodes into destination-oriented directed acyclic graphs (DODAGs) \cite{rfc6550}. DODAG Information Objects (DIOs) carry routing information, DODAG Information Solicitations (DISs) request it, and Destination Advertisement Objects (DAOs) advertise destination reachability. Albinali and Azzedin review attacks on these control messages \cite{albinali2024}. Previous work has studied repeated DIS solicitations \cite{rajasekar2022,vermaranga2020,medjek2021} and malicious DAO advertisements \cite{pu2018,limsd2024}. This study will examine a threat in which each compromised node performs both actions. Because the actions interfere with different routing functions, detection rate alone would provide an incomplete account of their effect.

We propose Adaptive Evidence Fusion for RPL (AEF-RPL), an extension of DIS-FDM. DIS-FDM uses Dempster--Shafer fusion to combine local observations of DIS traffic and preferred-parent changes \cite{prajisha2025}. The proposed extension adds a DAO-consistency check: when a node receives a DAO directly from its originator, it will compare the advertised transit parent with the observed relationship. AEF-RPL will integrate a tabular Q-learning decision layer, whose contribution will be evaluated against a learning-disabled variant. Both variants will use the same DAO observation so that the comparison can distinguish the effect of added evidence from the effect of learning.

Concurrent RPL attacks have been studied before. Garcia Ribera et al. ran Hello, DIS, and DAO floods together, with different nodes assigned to each attack \cite{garciaribera2022}. Other studies have examined complex attack combinations and joint-attack detection \cite{alsukayti2023,de2rarpl2024,abdulkareem2024}. These precedents motivate a more specific investigation: how a local detector responds when each malicious node executes DIS flooding and incorrect-parent DAO manipulation together, and how much DAO evidence and learning contribute separately. The proposed study will address this question through matched comparisons of the evidence and decision mechanisms. We use \emph{coordinated} in this limited same-node sense; the threat model will not require attackers to exchange plans through a separate protocol.

The proposed detector will operate locally, using a designated benign calibration period before combining evidence and applying selective DAO handling. Its learning layer will choose whether to accept or isolate a neighbor, subject to explicit decision rules. Earlier RPL defenses have also used Q-learning \cite{qsecrpl2023,itrpl2024}; its inclusion alone is therefore not a novelty claim. The evaluation will examine its contribution through false positives, alert delay, and communication performance, without assuming that adaptation necessarily improves these outcomes.

The study has three objectives: to assess an independent DIS-FDM reproduction, to develop the combined attack model and DAO-aware detector, and to compare their behavior under paired simulation conditions. Separating DAO evidence from the learning decision will support analysis of each component's contribution. The research significance lies in relating local evidence and mitigation decisions to detection and routing outcomes. The applied significance concerns the communication a defense may preserve and the risk of isolating honest nodes.

These objectives lead to five research questions:
\begin{enumerate}
\item[\emph{RQ1:}] To what extent does an independent Contiki-NG/Cooja implementation reproduce the network-performance trends reported for DIS-FDM under multicast DIS flooding?
\item[\emph{RQ2:}] How does concurrent DIS flooding and DAO manipulation affect packet delivery, delay, control overhead, and resilience compared with either attack alone?
\item[\emph{RQ3:}] To what extent does adding calibrated DAO-consistency evidence change detection performance and routing service compared with DIS-FDM and the selected DAO-focused defense?
\item[\emph{RQ4:}] How does local Q-learning change false positives, detection delay, and resilience compared with the same detector using a fixed decision rule?
\item[\emph{RQ5:}] How consistent are these outcomes across repeated seeds and static and mobile conditions, and what feasibility constraints arise when extending the evaluation from 30 to 500 and 1000 nodes?
\end{enumerate}

The planned contributions are:
\begin{enumerate}
\item A Contiki-NG/Cooja framework for examining multicast DIS flooding and transit-parent DAO manipulation from the same compromised nodes, with a local detector that combines observations of both behaviors under explicit first-hop consistency assumptions.
\item A paired evaluation of individual and combined attacks across attacker fractions and static and mobile networks. Repeated scenarios will compare detection, upward and downward delivery, delay, control traffic, and radio-time-based power estimates.
\item A controlled assessment of the Q-learning decision layer against fixed decisions using the same DAO mechanism. Attack-free mobility controls and component ablations will examine false isolation and the contribution of the main design choices.
\item A reproducibility package connecting scenario manifests, raw logs, firmware, configuration metadata, and analysis scripts. It will document assumptions and implementation differences from the source methods and support regeneration of the resulting tables and figures.
\end{enumerate}

The evaluation will focus initially on 30-node simulated RPL networks, with extensions to 500 and 1000 nodes subject to attack-free connectivity validation. The scope is the specified DIS/DAO behavior and local detection architecture; conclusions about other attack families, physical deployments, and alternative MAC protocols will require separate evidence. Detection quality, communication service, and false isolation will be assessed together rather than treating any single metric as sufficient evidence of effectiveness.

Section~\ref{sec:related} reviews the methods and evidence that motivate these questions, organized by control-message threat, detection mechanism, and decision placement.

\section{Literature Review}\label{sec:related}
We review control-plane defenses according to the information they inspect and the decisions they make from it. DIS-specific methods are considered first, followed by DAO and concurrent-attack detection, adaptive monitoring, and learning. Most of the application studies were published in 2021--2026. Earlier references supply the protocol specification and the foundations of methods used here. The purpose of this narrative review is to establish the closest technical comparisons; it does not claim complete coverage of RPL security.

\begin{table*}[t]
\caption{Positioning of AEF-RPL relative to selected RPL security methods.}
\label{tab:positioning}
\centering
\footnotesize
\setlength{\tabcolsep}{5pt}
\renewcommand{\arraystretch}{1.16}
\begin{tabularx}{\textwidth}{@{}p{0.14\textwidth}p{0.21\textwidth}p{0.25\textwidth}X@{}}
\toprule
Method & Threat or task & Main mechanism & Relationship to this study \\
\midrule
DIS-FDM \cite{prajisha2025} & Multicast DIS flooding & Distributed fusion of DIS and parent-change evidence & Methodological foundation; the proposed extension will add DAO consistency and evaluate learning separately. \\
Li-MSD \cite{limsd2024} & DAO insider attack & Lightweight blacklisting & DAO-focused comparator; attack semantics and implementation differences must be retained. \\
Garcia Ribera et al.\ \cite{garciaribera2022} & Multiple RPL attacks, including concurrent flooding & Hybrid placement and several detection rules & Prior concurrent multi-vector detection; the proposed combination will originate from each malicious node. \\
QSec-RPL \cite{qsecrpl2023} & Version-number attack under mobility & Q-learning-based detection & Prior learning-based RPL security for a different control-plane behavior. \\
iTRPL \cite{itrpl2024} & Trust-based routing decisions & Trust and multi-agent reinforcement learning & Different decision architecture and reward organization. \\
IMDRPL \cite{imdrpl2023} & DIO/DAO-related device security & Encrypted certificate attestation & Multi-message protection with different identity assumptions from local consistency checking. \\
FLSec-RPL / DSAT-IDS \cite{flsecrpl2024,dsatids2024} & DIO suppression / selective forwarding & Fuzzy validation / adaptive forwarding thresholds & Adaptive detection precedents; their observed signals differ from the proposed DIS/DAO observation. \\
CRRDS \cite{haqdad2026} & Decreased-rank attack & Deep reinforcement learning at the root & Contrasts centralized observation and resources with local tabular decisions. \\
AEF-RPL (proposed) & Concurrent DIS flood and incorrect-parent DAO & Local evidence fusion, calibration, DAO handling, tabular Q-learning & Planned paired assessment of added evidence and adaptation, including false-positive controls. \\
\bottomrule
\end{tabularx}
\end{table*}
\subsection{RPL Security and Operational Evaluation}
A control-message count and a routing inconsistency describe different aspects of an attack. RFC~6550 specifies how DIS, DIO, and DAO messages participate in routing \cite{rfc6550}; Albinali and Azzedin organize attacks and defenses across this control plane \cite{albinali2024}. The distinction matters when selecting DAO comparators. Repeated advertisements can be recognized through their rate, whereas an incorrect transit parent requires an observation of the advertised relationship. The planned comparison will retain these threat distinctions rather than group all DAO attacks under one detection task.

Detection and route recovery represent different evaluation targets. An attacker-level flag does not by itself establish that a damaged route has recovered. Similarly, a sent-versus-received delivery ratio excludes nodes that never transmit because they lack a usable route. The proposed evaluation will therefore examine application participation and network-wide service alongside detection and conditional delivery ratios. This distinction follows from the routing task and does not presume how any evaluated method will perform.

\subsection{DIS-Flooding Detection and Evidence Fusion}
The DIS literature establishes the traffic burden that motivates the baseline comparisons. Rajasekar and Rajkumar examine flooding impact \cite{rajasekar2022}; the mitigation studies of Verma and Ranga \cite{vermaranga2020} and Medjek et al.\ \cite{medjek2021} address DIS abuse in RPL networks. For the proposed combined threat, reducing that burden will be only part of the task. The same malicious node will also change the transit-parent information carried by its DAOs, so downward delivery will need to be examined alongside control traffic.

DIS-FDM provides the evidence-fusion foundation for the proposed AEF-RPL design \cite{prajisha2025}. The extension will add DAO-consistency observation and associated handling. Comparing AEF-RPL with DIS-FDM will therefore assess the extension as a whole. A separate learning-enabled versus learning-disabled comparison, retaining the same DAO mechanism, will be needed to isolate the decision layer's contribution.

A moving node may lose its preferred parent and select another without any adversary being present. Such changes complicate the use of parent instability as attack evidence. Mobile, attack-free controls will be included to measure false isolation under legitimate behavior. Detection rates obtained with attackers present cannot answer that question on their own.

\subsection{DAO Defenses and Concurrent Attacks}
The DAO-focused studies cover distinct behaviors. Pu examines DAO inconsistency \cite{pu2018}, while Li-MSD addresses a DAO insider attack \cite{limsd2024}. Li-MSD will serve as the DAO-focused comparator, while the proposed study will retain a specific threat definition: the originating attacker advertises an incorrect transit parent. AEF-RPL's consistency check will be restricted to the relationship observed when a DAO arrives directly from its originator. This observation point limits the scope of the proposed protection; it cannot be assumed to provide general resistance to DAO forgery.

The IDS of Garcia Ribera et al. has both a central component and distributed detectors \cite{garciaribera2022}. Its concurrent flooding test assigns Hello, DIS, and DAO attacks to different nodes. This provides a direct precedent for simultaneous attacks on several control-message types. The proposed scenario will instead assign both DIS and DAO actions to each attacker and evaluate local decisions. The comparison will examine effects on routing service as well as attacker flags.

Other relevant work includes Alsukayti and Alreshoodi's experiments with simple and complex routing attacks \cite{alsukayti2023}, DE2RA-RPL's treatment of resource-related attacks \cite{de2rarpl2024}, and the collaborative joint-attack IDS of Abdulkareem et al.\ \cite{abdulkareem2024}. Rouissat et al. study DIO inundation, which targets another control-message type \cite{rouissat2023}. These studies show that combining RPL attacks is an established research direction. The proposed contribution will concern the specified same-node DIS/DAO combination and a controlled assessment of the evidence and handling used against it.

\subsection{Adaptive and Multi-Message Detection}
Adaptation is broader than learning an action policy. FLSec-RPL uses fuzzy decisions with validation and blocking for DIO neighbor suppression and is evaluated under static and mobile conditions \cite{flsecrpl2024}. DSAT-IDS derives adaptive thresholds from group or subtree forwarding observations to detect selective forwarding \cite{dsatids2024}. These methods use evidence suited to their respective threats. Neither a forwarding-rate change nor a DIO-suppression indicator is the same observation as an incorrect DAO parent. This difference motivates evaluating AEF-RPL's additional evidence separately from the contribution of its learner.

Another multi-message approach is IMDRPL. Its DIO/DAO design uses encrypted certificate attestation and is evaluated with Contiki/Cooja \cite{imdrpl2023}. The authentication mechanism gives the comparison a different starting point from behavioral evidence fusion. The proposed AEF-RPL design will check a locally observed sender-parent relationship without a corresponding certificate mechanism. Its consistency judgment will therefore depend on explicit addressing and first-hop assumptions.

\subsection{Learning, Adaptation, and Detection Assumptions}
QSec-RPL uses Q-learning to detect version-number attacks in mobile RPL networks \cite{qsecrpl2023}. iTRPL combines trust with multi-agent reinforcement learning and divides responsibilities between nodes and the root \cite{itrpl2024}. Both provide precedents for learning-based RPL security, but their threats and decision structures differ from the proposed design. The planned learning comparison will examine changes to a local detector's accept/isolate decisions given the evidence available to AEF-RPL.

For decreased-rank detection, CRRDS uses a deep reinforcement learning agent at a resource-rich root and collects observations from child nodes \cite{haqdad2026}. The proposed AEF-RPL architecture will distribute tabular decisions across nodes. The placement changes both the available observations and computing resources, so a performance result for one arrangement does not validate the other. CRRDS is relevant as a precedent for root-based learning, rather than an equivalent implementation of the proposed local decision task.

The proposed decision layer will follow tabular Q-learning \cite{watkins1992}, with reward derived from a local anomaly indication rather than an independently verified attacker label. An incorrect indication could consequently influence later action values. The fixed-rule comparison will assess the effect of enabling learning, while attack-free mobile controls will examine false isolation during normal behavior. These experiments will not, by themselves, establish lifelong retention or recognition of previously unseen attack families.

A different evaluation task appears in the multiclass RPL dataset of Panda and Muthuraman \cite{panda2026}. Their supervised classifiers assign extracted observations to attack classes. AEF-RPL will instead select actions during network operation, so acceptance or isolation may change subsequent communication. A numerical comparison would require a common observation unit, dataset, and labeling protocol. This work helps position the proposed local detector alongside data-driven RPL detection, without treating the published classification scores as directly comparable performance evidence.

Related combinations of fusion and learning exist in vehicular reputation systems. Gyawali et al. aggregate feedback with Dempster--Shafer reasoning at edge servers and use deep reinforcement learning for reputation updates \cite{gyawali2021}. Jasti et al. describe a multi-level reputation system combining Dempster--Shafer and reinforcement learning \cite{jasti2024}. Combining the techniques is therefore not, by itself, the proposed novelty. The study will investigate the evidence available within RPL and the consequences of using it for local accept/isolate decisions under the specified concurrent threat.

\subsection{Comparative Synthesis and Study Positioning}
Table~\ref{tab:positioning} summarizes the closest methods and their relationship to this study. DIS-FDM supplies the fusion approach, Li-MSD provides a DAO-focused comparison, and the learning-based studies inform the adaptation question.

The review motivates three linked comparisons: individual versus combined attacks, DIS-focused versus DAO-aware detection, and learning-enabled versus fixed decisions using the same observations. The planned evaluation will include attack-free mobility controls and service measurements to distinguish false isolation from successful mitigation. These comparisons will test the proposed contribution without assuming that added evidence or learning must improve every outcome. The resulting evidence will be needed before claims of detection effectiveness, service preservation, or scalability can be made.

\begin{thebibliography}{99}
\bibitem{rfc6550} T. Winter, P. Thubert, A. Brandt, J. Hui, R. Kelsey, P. Levis, K. Pister, R. Struik, J. P. Vasseur, and R. Alexander, ``RPL: IPv6 Routing Protocol for Low-Power and Lossy Networks,'' Internet Engineering Task Force, RFC 6550, Mar. 2012, doi: \href{https://doi.org/10.17487/RFC6550}{10.17487/RFC6550}.
\bibitem{albinali2024} H. Albinali and F. Azzedin, ``Toward RPL attacks and mitigation taxonomy: Systematic literature review approach,'' \emph{IEEE Trans. Netw. Serv. Manag.}, vol. 21, no. 5, pp. 5215--5238, 2024, doi: \href{https://doi.org/10.1109/TNSM.2024.3386468}{10.1109/TNSM.2024.3386468}.
\bibitem{rajasekar2022} V. R. Rajasekar and S. Rajkumar, ``A study on impact of DIS flooding attack on RPL-based 6LowPAN network,'' \emph{Microprocessors and Microsystems}, vol. 94, Art. no. 104675, 2022, doi: \href{https://doi.org/10.1016/j.micpro.2022.104675}{10.1016/j.micpro.2022.104675}.
\bibitem{vermaranga2020} A. Verma and V. Ranga, ``Mitigation of DIS flooding attacks in RPL-based 6LoWPAN networks,'' \emph{Trans. Emerg. Telecommun. Technol.}, vol. 31, no. 2, Art. no. e3802, 2020, doi: \href{https://doi.org/10.1002/ett.3802}{10.1002/ett.3802}.
\bibitem{medjek2021} F. Medjek, D. Tandjaoui, N. Djedjig, and I. Romdhani, ``Multicast DIS attack mitigation in RPL-based IoT-LLNs,'' \emph{J. Inf. Secur. Appl.}, vol. 61, Art. no. 102939, 2021, doi: \href{https://doi.org/10.1016/j.jisa.2021.102939}{10.1016/j.jisa.2021.102939}.
\bibitem{pu2018} C. Pu, ``Mitigating DAO inconsistency attack in RPL-based low power and lossy networks,'' in \emph{Proc. IEEE 8th Annu. Comput. Commun. Workshop Conf. (CCWC)}, 2018, pp. 570--574, doi: \href{https://doi.org/10.1109/CCWC.2018.8301614}{10.1109/CCWC.2018.8301614}.
\bibitem{limsd2024} A. Verma, S. K. Verma, A. C. Pandey, J. Grover, and G. Sharma, ``Li-MSD: A lightweight mitigation solution for DAO insider attack in RPL-based IoT,'' \emph{Future Gener. Comput. Syst.}, vol. 159, pp. 327--339, 2024, doi: \href{https://doi.org/10.1016/j.future.2024.05.032}{10.1016/j.future.2024.05.032}.
\bibitem{prajisha2025} C. Prajisha and A. R. Vasudevan, ``Detection and mitigation of multicast DIS flooding attacks in RPL-based IoT networks,'' \emph{Ad Hoc Networks}, vol. 179, Art. no. 104002, 2025, doi: \href{https://doi.org/10.1016/j.adhoc.2025.104002}{10.1016/j.adhoc.2025.104002}.
\bibitem{garciaribera2022} E. Garcia Ribera, B. Martinez Alvarez, C. Samuel, P. P. Ioulianou, and V. G. Vassilakis, ``An intrusion detection system for RPL-based IoT networks,'' \emph{Electronics}, vol. 11, no. 23, Art. no. 4041, 2022, doi: \href{https://doi.org/10.3390/electronics11234041}{10.3390/electronics11234041}.
\bibitem{alsukayti2023} I. S. Alsukayti and M. Alreshoodi, ``RPL-based IoT networks under simple and complex routing security attacks: An experimental study,'' \emph{Applied Sciences}, vol. 13, no. 8, Art. no. 4878, 2023, doi: \href{https://doi.org/10.3390/app13084878}{10.3390/app13084878}.
\bibitem{de2rarpl2024} D. Paganraj and M. Chelliah, ``DE2RA-RPL: Detection and elimination of resource-related attacks in IoT RPL-based protocol,'' \emph{J. Supercomput.}, vol. 80, no. 15, pp. 22397--22427, 2024, doi: \href{https://doi.org/10.1007/s11227-024-06256-w}{10.1007/s11227-024-06256-w}.
\bibitem{abdulkareem2024} O. A. Abdulkareem, R. K. Kontham, and F. E. Mahmood, ``Collaborative intrusion detection system to identify joint attacks in routing protocol for low-power and lossy networks routing protocol on the Internet of Everything,'' \emph{Mesopotamian J. CyberSecurity}, vol. 4, no. 3, pp. 251--277, 2024, doi: \href{https://doi.org/10.58496/MJCS/2024/026}{10.58496/MJCS/2024/026}.
\bibitem{qsecrpl2023} G. Sharma, J. Grover, and A. Verma, ``QSec-RPL: Detection of version number attacks in RPL based mobile IoT using Q-learning,'' \emph{Ad Hoc Networks}, vol. 142, Art. no. 103118, 2023, doi: \href{https://doi.org/10.1016/j.adhoc.2023.103118}{10.1016/j.adhoc.2023.103118}.
\bibitem{itrpl2024} D. Dey and N. Ghosh, ``iTRPL: An intelligent and trusted RPL protocol based on multi-agent reinforcement learning,'' \emph{Ad Hoc Networks}, vol. 163, Art. no. 103586, 2024, doi: \href{https://doi.org/10.1016/j.adhoc.2024.103586}{10.1016/j.adhoc.2024.103586}.
\bibitem{imdrpl2023} P. Deepavathi and C. Mala, ``IMDRPL: Identifying and eliminating malicious devices using DIO and DAO ICMP control messages in RPL-based protocol,'' \emph{Peer-to-Peer Netw. Appl.}, vol. 16, no. 5, pp. 2380--2398, 2023, doi: \href{https://doi.org/10.1007/s12083-023-01539-0}{10.1007/s12083-023-01539-0}.
\bibitem{flsecrpl2024} C. Kim, C. So-In, Y. Kongsorot, and P. Aimtongkham, ``FLSec-RPL: A fuzzy logic-based intrusion detection scheme for securing RPL-based IoT networks against DIO neighbor suppression attacks,'' \emph{Cybersecurity}, vol. 7, Art. no. 27, 2024, doi: \href{https://doi.org/10.1186/s42400-024-00223-x}{10.1186/s42400-024-00223-x}.
\bibitem{dsatids2024} V. R. Rajasekar and S. Rajkumar, ``DSAT-IDS: Dissimilarity and adaptive threshold-based intrusion detection system to mitigate selective forwarding attack in the RPL-based 6LoWPAN,'' \emph{Cluster Comput.}, vol. 27, no. 9, pp. 12029--12068, 2024, doi: \href{https://doi.org/10.1007/s10586-024-04564-7}{10.1007/s10586-024-04564-7}.
\bibitem{haqdad2026} M. Haqdad, M. Fayaz, P. Khan, F. Ali, T. H. H. Aldhyani, A. K. Bashir, and D. Kwak, ``Intent-based networking with deep reinforcement learning for detecting decreased rank attacks in low-power and lossy IoT networks,'' \emph{IEEE Internet Things J.}, vol. 13, no. 9, pp. 17802--17814, 2026, doi: \href{https://doi.org/10.1109/JIOT.2025.3596097}{10.1109/JIOT.2025.3596097}.
\bibitem{rouissat2023} M. Rouissat, M. Belkheir, I. S. Alsukayti, and A. Mokaddem, ``A lightweight mitigation approach against a new inundation attack in RPL-based IoT networks,'' \emph{Applied Sciences}, vol. 13, no. 18, Art. no. 10366, 2023, doi: \href{https://doi.org/10.3390/app131810366}{10.3390/app131810366}.
\bibitem{watkins1992} C. J. C. H. Watkins and P. Dayan, ``Q-learning,'' \emph{Machine Learning}, vol. 8, nos. 3--4, pp. 279--292, 1992, doi: \href{https://doi.org/10.1007/BF00992698}{10.1007/BF00992698}.
\bibitem{panda2026} N. Panda and S. Muthuraman, ``A multiclass machine learning framework for detecting routing attacks in RPL-based IoT networks using a novel simulation-driven dataset,'' \emph{Future Internet}, vol. 18, no. 1, Art. no. 35, 2026, doi: \href{https://doi.org/10.3390/fi18010035}{10.3390/fi18010035}.
\bibitem{gyawali2021} S. Gyawali, Y. Qian, and R. Q. Hu, ``Deep reinforcement learning based dynamic reputation policy in 5G based vehicular communication networks,'' \emph{IEEE Trans. Veh. Technol.}, vol. 70, no. 6, pp. 6136--6146, 2021, doi: \href{https://doi.org/10.1109/TVT.2021.3079379}{10.1109/TVT.2021.3079379}.
\bibitem{jasti2024} P. C. Jasti, H. Griffith, and H. Rathore, ``A multi-level Dempster-Shafer and reinforcement learning-based reputation system for connected vehicle security,'' in \emph{Proc. IEEE 21st Consum. Commun. Netw. Conf. (CCNC)}, 2024, pp. 714--720, doi: \href{https://doi.org/10.1109/CCNC51664.2024.10454856}{10.1109/CCNC51664.2024.10454856}.
\end{thebibliography}
\end{document}
